%==============================================================================
% Section 8 -- Optimisation: Densification, Pruning, and Opacity
%==============================================================================
\section{Optimisation Mechanisms: Densification, Pruning, and Opacity}
\label{sec:optimization}

\subsection{Training Objective}

All parameters are optimised end-to-end by minimising:
\begin{equation}
  \boxed{
    \mathcal{L} = (1 - \lambda)\,\mathcal{L}_1 + \lambda\,\mathcal{L}_{\text{D-SSIM}}
  }, \quad \lambda = 0.2
\end{equation}
$\mathcal{L}_1$ ensures per-pixel accuracy; $\mathcal{L}_{\text{D-SSIM}}$
preserves structural detail by penalising blurry renderings.

%------------------------------------------------------------------------------
\subsection{Optimiser}

Adam with per-group learning rates (position: $1.6\times10^{-4}$ decaying to
$1.6\times10^{-6}$; scale: $5\times10^{-3}$; rotation: $10^{-3}$; opacity:
$5\times10^{-2}$; SH degree 0: $2.5\times10^{-3}$; SH 1--3:
$2.5\times10^{-4}$).

%------------------------------------------------------------------------------
\subsection{Adaptive Density Control}
\label{ssec:density}

Every 100 iterations (from iter.\ 500 to 15\,000), the mean position gradient
\begin{equation}
  \bar{\nabla}_i = \tfrac{1}{K}\sum_{k=1}^{K}\|\nabla_{\mean'_i}\mathcal{L}\|
\end{equation}
drives two operations when $\bar{\nabla}_i > \tau_{\text{grad}} = 2\times10^{-4}$:

\begin{itemize}
  \item \textbf{Clone} (small Gaussian, $\max s_k < \tau_{\text{size}}$): create
        an exact copy; both are optimised independently. The region needs more
        primitives.
  \item \textbf{Split} (large Gaussian, $\max s_k \geq \tau_{\text{size}}$):
        replace with two Gaussians at scale $s_k/1.6$, positions sampled from
        $\mathcal{N}(\mean_i, \cov_i)$. The primitive covers too large an area.
\end{itemize}

%------------------------------------------------------------------------------
\subsection{Pruning}

At the same frequency, Gaussians are removed if: (1) opacity
$o_i < \epsilon_\alpha = 0.005$, (2) world-space size exceeds $\tau_w$, or
(3) projected 2D radius exceeds $\tau_v$ (removes ``floaters'').

%------------------------------------------------------------------------------
\subsection{Periodic Opacity Reset}

Every 3000 iterations all opacities are reset to $\alpha_{\min} \approx 0.01$.
Gaussians that contribute to the scene recover high opacity quickly; redundant
ones do not and are pruned at the next pruning step. This prevents ghost
Gaussians from accumulating throughout training.

%------------------------------------------------------------------------------
\subsection{Training Schedule and Convergence}

\begin{algorithm}[H]
  \caption{3DGS Training Schedule (30\,000 iterations)}
  \label{alg:training}
  \begin{algorithmic}[1]
    \Require{Images $\{I_k\}$, SfM point cloud $\mathcal{P}$}
    \State Initialise small isotropic Gaussians from $\mathcal{P}$
    \For{$t = 1, \ldots, 30000$}
      \State Render $\hat{I}$; compute $\mathcal{L}$; Adam step
      \If{$500 \leq t \leq 15000$ and $t \bmod 100 = 0$}
        \State Clone/split high-gradient Gaussians; prune low-opacity/large ones
      \EndIf
      \If{$t \bmod 3000 = 0$} Reset opacities to $\alpha_{\min}$ \EndIf
      \If{$t \bmod 1000 = 0$ and $t \leq 3000$} Unlock next SH degree \EndIf
    \EndFor
  \end{algorithmic}
\end{algorithm}

A typical run starts with $\sim10^4$--$10^5$ Gaussians, peaks at
$\sim3\times10^6$--$6\times10^6$, and stabilises at
$\sim10^6$--$4\times10^6$. Training takes 20--45 min on an RTX 3090/4090;
at convergence rendering exceeds 90 fps at $1920\times1080$.
